\section*{Appendici}
\addcontentsline{toc}{section}{Appendici}

\subsection*{Appendice A  Materiale di riferimento}
Il piano si basa su: 

\paragraph{Documenti di progetto}
\begin{itemize}
    \item \textbf{D0  Proposta di progetto}: composizione del team e curriculum dei membri.
    \item \textbf{D1  Analisi del contesto e dei requisiti}: stakeholder, risultati del questionario, elenco dei requisiti RF01RF32 con priorità, analisi dei competitor e dei mockup, strumenti di sviluppo.
\end{itemize}

\paragraph{Materiale del committente e degli utenti}
\begin{itemize}
    \item Indicazioni del docente sul progetto e calendario dei deliverable D1D7.
    \item Incontro online su Zoom, organizzato dal docente, con i responsabili del Laboratorio in Campo di Scienze Ambientali.
    \item Questionario (Google Form) a studenti e docenti del laboratorio, somministrato dal 13 al 15 ottobre 2026 (7 risposte).
\end{itemize}

\paragraph{Mock-up dell'hackathon di settembre}
\begin{itemize}
    \item \textbf{BallMap}: 14 schermate (accesso, mappa, rilevazione, storico, revisione e permessi).
    \item \textbf{Field Lab}: 23 schermate (dashboard, botanica, chimica dell'acqua, gruppo e chat, report, profilo, sincronizzazione e SOS).
\end{itemize}
L'analisi dettagliata di entrambi è nella Sezione~4.2 del D1.

\paragraph{Applicazioni analizzate}
QField (per QGIS), Rock Identifier e KoboToolbox/KoboCollect, confrontate nella Sezione~4.1 del D1.


